% Pro M. Caelio (pro-caelio) --- English translation
% Translated by Claude (Anthropic) from the Latin (Perseus).
% Style guide: STYLE.md.

\ciceroSection{1} If anyone, members of the jury, were now perchance present, ignorant of our laws, our courts, our custom, he would surely wonder what so great atrocity there is in this case, that on festive days and during the public games, with every other forensic business broken off, this one trial alone is being held --- and he would not doubt that the defendant is charged with so great a crime that, if it were neglected, the state could not stand. The same man, when he heard that there is a law which orders daily inquiry into seditious and ruined citizens who in arms have besieged the senate, have brought violence on magistrates, have made war on the commonwealth, would not condemn the law, but would ask what crime was being moved in the trial. When he heard that no crime, no audacity, no violence is being called to trial, but that a young man of brilliant talent, industry, and popularity is being prosecuted by the son of a man whom he himself is bringing, and has brought, into court, while the resources of a courtesan are being used to assault him --- he would not censure the dutifulness of Atratinus himself, would judge that a woman's lust must be checked, and would think you, gentlemen, hard-driven, since not even in the common leisure of the city are you allowed to be at leisure.

\ciceroSection{2} For if you are willing to attend carefully and to judge truly of this entire case, you will conclude, members of the jury, that no one would have come down to this prosecution to whom it had been left to choose, nor, when he had come down, would have had any hope, unless he were leaning on someone's intolerable lust and excessively bitter hatred. But Atratinus, a most cultured and excellent young man, my close connection, I forgive, since he has the excuse either of dutifulness, or of necessity, or of his age. If he wished to prosecute, I attribute it to dutifulness; if he was ordered, to necessity; if he hoped for something, to youth. To the rest, not only is nothing to be forgiven, but they are to be sharply resisted.

\ciceroSection{3} And to me, members of the jury, this seems the entrance most suited to the defence of the youth of M. Caelius: that I should first reply to the things that the prosecutors said for the disfiguring of this case and the dragging down and stripping of his standing. His father has been thrown at him in various ways --- said either to be too little distinguished himself, or to have been treated too little dutifully by his son. As to standing: M. Caelius's father easily and silently answers for himself, even without my speech, before those who know him and his elders. To those, however, who, on account of his age (since he has now for some time spent less time in the Forum and with us), have not got an equal acquaintance with him, let them hold this: that whatever standing there can be in a Roman knight --- and that can certainly be the very greatest --- has always been held the highest in M. Caelius, and is held so today, not only by his own circle but by everyone to whom for any cause he could be known.

\ciceroSection{4} Moreover, that being the son of a Roman knight should be set down as a charge by prosecutors --- this neither these men ought to have allowed in their pleading, nor we in our defence. As to what you said about his dutifulness, this estimate is indeed ours, but the judgement is certainly the parent's. What we hold, you will hear from those under oath; what the parents feel, the tears of his mother and her unbelievable grief, the squalor of his father, and this present sorrow that you see, and his mourning, declare.

\ciceroSection{5} As to the charge that the young man is not approved by his fellow townsmen --- the people of \dots\ have never given greater honours to anyone present, members of the jury, than to M. Caelius in his absence: him in his absence they co-opted into the most ample order, and they conferred on him, when he was not even seeking them, things which they had refused to many seekers. The same men have now sent the most select men of our order and Roman knights with a delegation to this trial and with the most weighty and ornate testimonial. I seem to myself to have laid the foundations of my defence, the firmest of all if they rest on the judgement of his own people. Nor indeed could his youth be sufficiently commended to you if it were displeasing not only to a parent of such quality, but also to a township so distinguished and so weighty. As for myself,

\ciceroSection{6} to come back to me, it is from these same fountains that I have flowed forth into the public reputation of men, and this work of mine in the Forum and the order of my life has spread out into men's estimation a little more widely on the strength of the commendation and judgement of my own. As to what was thrown at him on the score of his chastity, and what was kept up by all his prosecutors not by way of charges but of voices and abuses --- this M. Caelius will never take so bitterly as to repent of not having been born ill-favoured. For these are the abuses widespread against all whose youth had a free-born form and look. But it is one thing to abuse, another to prosecute. Prosecution requires a charge that it may define the matter, mark out the person, prove by argument, confirm by witness; abuse has no aim except contumely, which, if it is thrown more saucily, is called railing, if more wittily, urbanity.

\ciceroSection{7} Which part of the prosecution I was indeed astonished at, and took it ill, that it had been chiefly entrusted to Atratinus. For neither was it fitting, nor did his age call for it, nor (as you could see) did the modesty of an excellent young man permit him to engage in such a speech. I could have wished some of you sturdier men had taken on this office of abuse: somewhat more freely and more strongly and more after our manner we should refute that licence of slander. With you, Atratinus, I shall deal more gently, both because your modesty restrains my speech, and because I owe it to my service to you and to your father to protect him. There is, however, this I want you to be admonished of:

\ciceroSection{8} first, that what you are, all may judge you to be --- that as far as you stand off from foulness in your conduct, so far you separate yourself from licence in your words; next, that you do not say of another things at which, if they were falsely returned upon you, you would blush. For who is there to whom that road does not lie open? who is there who could not, as petulantly as he wished, abuse this age and this standing --- even if without the least suspicion, yet not without an argument? But the fault of those parts of the case is on the men who wished you to plead them; the praise of your own modesty is that we saw you saying these things unwillingly, of your talent, that you said them ornately and elegantly.

\ciceroSection{9} But to that whole speech the defence is brief. For so far as the age of M. Caelius could give a foothold to that suspicion, it was first defended by his own modesty, and then also by his father's care and discipline. From the time the father gave him the toga of manhood --- I shall say nothing here of myself; let it be only as much as you yourselves estimate; this I shall say, that he was led from his father straight to me --- no one ever saw this M. Caelius in that bloom of his age except either with his father or with me, or being trained in the most chaste house of M. Crassus in the most honourable arts.

\ciceroSection{10} As to the familiarity with Catiline that has been thrown at Caelius, that ought to be far removed from such a suspicion. For you know that, when Caelius was already a young man, Catiline stood for the consulship with me. If he ever went over to him, or ever fell off from me --- although many fine young men were attached to that worthless and wicked man --- then let Caelius be reckoned to have been too closely Catiline's familiar. ``But afterwards we know and we saw that this man was even among that man's friends.'' Who denies it? But this period of life, which itself is by its own nature weak, and harassed besides by the lust of others, this it is that I am here defending. He was constantly with me when I was praetor; he did not know Catiline. Africa Catiline was at that time holding as praetor. Then followed the year in which Catiline pleaded his case on the charge of extortion. Caelius was with me; to that man he never came even as a supporter. Then came the year in which I stood for the consulship; Catiline was standing with me. Caelius never went near him; from me he never withdrew.

\ciceroSection{11} For so many years, then, he had passed his life in the Forum without suspicion, without ill repute, when he attached himself to Catiline standing a second time. To what limit, then, do you suppose his youth ought to have been kept under guard? In our time at least one year used to be set down for keeping the arm under the toga; for using exercise and the games of the Campus we were in the tunic; and the same was the rule of the camp and of the soldier's life, if we had at once begun to serve. At that age, unless a man defended himself by his own gravity, his chastity, both his domestic discipline and a certain natural goodness besides, however he had been kept watch over by his own people, he could not still escape genuine ill repute. But of him who had kept those first beginnings of his age sound and untouched, of his fame and his chastity, when now he had braced himself and was a man among men, no one spoke.

\ciceroSection{12} ``But Caelius attached himself to Catiline, when he had already been some years in the Forum.'' And many of every order and every age did the same. For Catiline, as I think you remember, possessed --- not painted out, but lightly sketched in --- very many likenesses of the greatest virtues. He used many wicked men, and yet he pretended to be devoted to the best of men. There were in him many enticements to lust; there were also certain spurs to industry and labour. The vices of lust were blazing in him; there flourished too the studies of war. Nor do I think there ever was such a monster on earth, blended together as he was out of contrary and various and warring natural pursuits and desires.

\ciceroSection{13} Who, at one time, was more agreeable to the more distinguished men, who more closely tied to the more disgraceful? who was sometimes a citizen of the better party, who a more dreadful enemy to this state? who in pleasures more polluted, who in labours more enduring? who in rapacity more grasping, who in lavishness more pouring out? Those qualities indeed, gentlemen, in that man were astonishing: to embrace many men in friendship, to retain them by attentiveness, to share with all what he had, to serve every man's needs with his money, his influence, his bodily labour, even with crime, if it were needed, and audacity; to turn his nature and govern it according to occasion, to twist and bend it this way and that --- with the strict to live strictly, with the relaxed pleasantly, with the elderly seriously, with the youthful affably, with criminals audaciously, with the libertines in luxury.

\ciceroSection{14} With this so varied and manifold a nature, when he had gathered every wicked and audacious man out of every land, so too he held many brave and good men by a certain semblance of a counterfeit virtue. Nor would that abandoned onset of destroying this empire ever have come from him, if so vast a savagery of so many vices had not rested on certain roots of complaisance and forbearance. Therefore let this charge, members of the jury, be spat out, and let no charge of familiarity with Catiline cling. For it is shared with many, even with some good men. He almost deceived even me --- me, I say --- when he seemed to me both a good citizen and eager for every best man and a steadfast and faithful friend; and his deeds I caught with my eyes before my opinion, with my hands before my suspicion. Whose great packs of friends, if Caelius too was among them, it is rather for him to be sorry to have made a mistake, just as I sometimes regret my own mistake about that same man, than for him to fear the charge of that friendship.

\ciceroSection{15} And so your speech has slid down from the abuses of unchastity to the odium of the conspiracy. For you laid down --- yet that hesitantly and curtly --- that this man was a partner of the conspiracy on account of his familiarity with Catiline; in which not only the charge would not stick, but the speech of an eloquent young man scarcely held together. For what so great madness was there in Caelius, what so great a wound either in his character and nature, or in his property and fortune? Where, indeed, was Caelius's name ever heard in that suspicion? I am speaking too much about a matter not in the least doubtful; this, however, I do say. Not only if he had been a partner in the conspiracy, but unless he had been a most bitter enemy of that crime, he would never have wished, by an accusation of conspiracy, to commend his own youth in particular.
